\chapter{Reproducibility Details}
\label{app:reproducibility}

This appendix collects the hyperparameters, reward-profile definitions, and supplementary results that
support the main text. Values still to be finalized from the experiment runs are marked.

\section{Single-Path Selection: Reward Profiles}
\label{app:p1profiles}

The conditional path-selection agent is trained and evaluated across the reward-weight vectors
$\mathbf{w} = (w_1, w_2, w_3, w_4, w_{\text{probe}})$ of \cref{eq:p1reward}, spanning distinct application
intents. \Cref{tab:app:profiles} lists the six profiles used, together with the label each carries in
\cref{ch:single_path}.

\begin{table}[h]
  \centering
  \caption[The six reward profiles of the single-path study]{The six reward profiles (application intents) used in the single-path selection study
  (\cref{ch:single_path}). $w_1$ weights goodput, $w_2$ trust, $w_3$ loss and $w_4$ delay within the
  trust term, and $w_{\text{probe}}$ the probing penalty. All six are used in training on a stratified
  schedule; the four marked with a chapter label are also used in the evaluation of
  \cref{sec:sp:eval}.}
  \label{tab:app:profiles}
  \begin{tabular}{llccccc}
    \toprule
    Profile & Label in \cref{ch:single_path} & $w_1$ & $w_2$ & $w_3$ & $w_4$ & $w_{\text{probe}}$ \\
    \midrule
    \texttt{bandwidth\_max}    & Throughput  & 1.00 & 0.00 & 0.05 & 0.05 & 0.05 \\
    \texttt{delay\_averse}     & Low-Latency & 0.10 & 0.90 & 0.15 & 1.00 & 0.05 \\
    \texttt{loss\_averse}      & Low-Loss    & 0.05 & 0.95 & 1.00 & 0.15 & 0.05 \\
    \texttt{balanced\_extreme} & Balanced    & 0.55 & 0.45 & 0.55 & 0.55 & 0.05 \\
    \texttt{probe\_minimal}    & --          & 0.70 & 0.30 & 0.50 & 0.50 & 0.001 \\
    \texttt{probe\_averse}     & --          & 0.70 & 0.30 & 0.50 & 0.50 & 0.35 \\
    \bottomrule
  \end{tabular}
\end{table}

\section{Single-Path Selection: Agent Hyperparameters}
\label{app:p1hyper}

\Cref{tab:app:p1hyper} lists the hyperparameters shared by the flat, scoring, and conditional agents of
\cref{sec:p1impl:agents}; the flat agent differs only in training for 533~episodes rather than 666.

\begin{table}[h]
  \centering
  \caption[Hyperparameters for the single-path agents]{Hyperparameters for the single-path selection agents (\cref{ch:single_path}).}
  \label{tab:app:p1hyper}
  \begin{tabular}{ll}
    \toprule
    Parameter & Value \\
    \midrule
    Hidden layers / width          & 2 / 128 \\
    Learning rate (Adam)           & $10^{-3}$ \\
    Discount $\gamma$              & 0.95 \\
    Batch size                     & 32 \\
    Replay capacity                & $10^4$ (minimum 500 before learning) \\
    Target update                  & soft (Polyak), $\tau = 0.05$, every 100 gradient steps \\
    Prioritized replay             & $\alpha = 0.6$, $\beta$ annealed $0.4 \rightarrow 1.0$ \\
    Exploration                    & $\varepsilon$-greedy, $1.0 \rightarrow 0.01$, decay 0.995/episode \\
    Dropout                        & 0.1 \\
    Gradient clipping              & 1.0 \\
    Episodes $\times$ steps        & $666 \times 24$ (flat: $533 \times 24$) \\
    Gradient step frequency        & every 4 environment steps \\
    Invalid-action mask value      & $-10^{9}$ \\
    Random seeds                   & 5 (all agents) \\
    Wall-clock per agent           & 1--3\,min (22-core CPU, no GPU) \\
    \bottomrule
  \end{tabular}
\end{table}

Two caveats on this table. The $\varepsilon$ schedule is stated as its configured range, but a run is too
short to traverse it: at 0.995 decay per episode, $\varepsilon$ reaches only 0.036 after 666~episodes and
0.069 after 533, so the 0.01~floor is never attained and all reported agents retain a few percent of
exploration at the end of training. Seeding fixes network initialization, prioritized-replay sampling,
and $\varepsilon$-greedy exploration; the environment's pair, hour, and profile streams are deliberately
left fixed across seeds, so every seed sees the identical stream of training contexts and only the
learning process varies.

\section{Multipath System: QoE Weights and SAC Hyperparameters}
\label{app:p2hyper}

The QoE reward of \cref{eq:p2qoe} uses base weights $(a,b,c,d,e) = (1, 0.5, 0.5, 1, 0.25)$, overridden
per intent by \cref{tab:p2profiles}; the application agent acts every \SI{1}{\second} and the transport
agent every frame (${\approx}\SI{33}{\milli\second}$ at 30\,fps) over a bitrate range of
\SIrange{300}{6000}{\kilo\bit\per\second}. \Cref{tab:app:p2hyper,tab:app:p2topo} give the learner and
scenario parameters.

\begin{table}[h]
  \centering
  \caption[Soft actor-critic hyperparameters, multipath system]{Soft actor-critic hyperparameters shared by the application and transport agents
  (\cref{sec:p2impl:agents}). Source: \texttt{rl-mpquic/configs/dynamic.yaml}.}
  \label{tab:app:p2hyper}
  \begin{tabular}{ll}
    \toprule
    Parameter & Value \\
    \midrule
    Hidden layers / width       & 2 / 256 \\
    Learning rate (Adam)        & $3\times10^{-4}$ \\
    Discount $\gamma$           & 0.99 \\
    Target update               & soft (Polyak), $\tau = 0.005$ \\
    Batch size                  & 256 \\
    Replay capacity             & $2\times10^{5}$ \\
    Random warm-up / first update & 1000 / 1000 environment steps \\
    Temperature $\alpha$        & auto-tuned, target entropy $-d_{\mathrm{act}}$ \\
    Critic LayerNorm (DroQ)     & enabled \\
    Updates per environment step & 2 \\
    Prioritized replay          & enabled \\
    Episodes $\times$ length    & $100 \times \SI{30}{\second}$ (900~frames) \\
    Wall-clock per run          & 4--4.5\,h on the mock, 20--25\,h on NS-3 (16-core CPU, no GPU) \\
    \bottomrule
  \end{tabular}
\end{table}

\begin{table}[h]
  \centering
  \caption[The dynamic six-path scenario]{The dynamic six-path scenario (\cref{sec:p2eval:setup}): per-path bottleneck rate,
  propagation delay, and mean cross-traffic as a fraction of that rate. No path can carry full-quality
  video alone. Paths~5 and~6 form a shared-bottleneck group that degrades together. Source:
  \texttt{rl-mpquic/configs/dynamic.yaml}.}
  \label{tab:app:p2topo}
  \begin{tabular}{lccl}
    \toprule
    Path & Rate & Delay & Character \\
    \midrule
    1 & \SI{3.0}{\mega\bit\per\second} & \SI{10}{\milli\second} & clean, low RTT (cross 0.40) \\
    2 & \SI{2.5}{\mega\bit\per\second} & \SI{15}{\milli\second} & clean (cross 0.45) \\
    3 & \SI{2.5}{\mega\bit\per\second} & \SI{40}{\milli\second} & latency trap (cross 0.30) \\
    4 & \SI{2.0}{\mega\bit\per\second} & \SI{20}{\milli\second} & congested (cross 0.55) \\
    5 & \SI{3.0}{\mega\bit\per\second} & \SI{25}{\milli\second} & shared bottleneck A (cross 0.50) \\
    6 & \SI{2.5}{\mega\bit\per\second} & \SI{30}{\milli\second} & shared bottleneck B (cross 0.50) \\
    \bottomrule
  \end{tabular}
\end{table}

The non-stationary mechanisms are hazards converted per frame from per-second rates. Path churn takes a
live path down at \SI{0.08}{\per\second} and revives a dead one at \SI{0.30}{\per\second}, with at least
three of six paths live at all times. Regime shifts resample a path's capacity multiplier uniformly in
$[0.55, 1.25]$ at \SI{0.20}{\per\second}. Congestion bursts drop capacity to 45\% for
\SI{0.5}{\second} at \SI{0.10}{\per\second}. The correlated-failure group $\{5, 6\}$ drops to 50\% for
\SI{1.2}{\second} at \SI{0.05}{\per\second}.

\paragraph{Intent conditioning.} The conditioned agents of \cref{sec:p2eval:conditioning} append a
five-component intent vector to the observation: the delay, jitter, and loss weights $(b, c, d/2)$ of
\cref{eq:p2qoe}, the frame deadline, and the latency normalizer, the last two scaled by
\SI{800}{\milli\second}. The existing observation features are left untouched and the normalizer is
disclosed as an input rather than fixed across intents, which is what keeps the per-intent specialists
of \cref{tab:p2transfer} a valid baseline against the conditioned policy --- and is also the third of
the three premise violations reported in \cref{sec:p2design:cancel}, since that normalizer divides a
per-path feature. Two conditioning modes are compared: \texttt{global}, which supplies the vector to the
global context slot alone, and \texttt{perpath}, which additionally broadcasts it to every per-path
feature row, adding $5 \times 6 = 30$ input dimensions. Training randomizes the intent per episode,
stratified over the three anchors of \cref{tab:p2profiles} plus a continuously sampled interior; the
randomization is training-only and evaluation never resamples. Both arms are trained for six seeds at the
budget of \cref{tab:app:p2hyper}. Source: \texttt{rl-mpquic/configs/dynamic\_intent\{,\_glob\}.yaml}.

A third arm, \cref{tab:p2fixednorm}'s, is identical to the \texttt{perpath} one apart from
\texttt{intent.obs\_norm: fixed}, which pins the latency and jitter normalizers to the G.114 and Y.1541
constants (\SI{200}{\milli\second} and \SI{50}{\milli\second}) instead of taking them from the intent
bundle, and so removes the second, unintended intent channel into the per-path rows. Because the
normalizer feeds only the observation scaling --- never the reward, the deadline, or the dynamics ---
the two arms are graded on environments that differ in no other respect, each under the normalizer it
trained with, which is what makes the comparison paired. It is trained for three seeds at the same
budget; three seeds concurrent took 3\,h\,07 wall-clock. Source:
\texttt{rl-mpquic/configs/dynamic\_intent\_fixednorm.yaml}.

\paragraph{The NS-3-native arm.} The three packet-level-native seeds of \cref{sec:p2eval:ns3native} use
the same \texttt{configs/dynamic.yaml} and the same budget as every other arm --- 100 episodes of
\SI{30}{\second}, six paths, TCP subflows, all four non-stationary mechanisms enabled --- differing only
in \texttt{backend: ns3}. They are trained at seeds 1, 2, and 3, strictly sequentially for the
shared-memory reason given in \cref{sec:p2eval:ns3native}, and completed in 24.6, 25.0, and 20.1\,h. The
three gradings are one \texttt{ab\_eval} invocation each over the same 3~checkpoints $\times$ 3
evaluation seeds (1000, 2000, 3000) $\times$ 5 episodes: the arm on NS-3, the same checkpoints on the
mock, and a paired two-arm comparison against the three mock-trained \texttt{dyn-s\{1,2,3\}}
checkpoints on NS-3. Source: \texttt{rl-mpquic/runs/ns3train-s\{1,2,3\}/} and
\texttt{runs/ab-d5-\{ns3native,onmock,compare\}/}.

\paragraph{The two static scenarios.} \Cref{sec:p2eval:regimes} contrasts two stationary
configurations, and they are not interchangeable, so both are stated here. The \emph{asymmetric} one
has three paths at 8, 4, and \SI{2}{\mega\bit\per\second}, a four-fold spread: an even split places a
third of the load on a path holding an eighth of the capacity, so the weakest path is overdriven
continuously and the app-only ablation collapses even though nothing in the network moves. The
\emph{near-equal} one has four paths of comparable rate, so an even split fits inside every path's
residual capacity and app-only comes within 2\% of the full system. Neither enables any of the four
non-stationary mechanisms. The pair is what isolates headroom heterogeneity from motion: the network
is stationary in both, and only the spread of capacity across paths differs.

\section{Supplementary Results}
\label{app:extra}

\begin{figure}[h]
  \centering
  \includegraphics[width=\linewidth]{p1app_training_curves.png}
  \caption[Training curves for the single-path agents]{Training curves for the single-path agents of \cref{sec:p1impl:agents}: episode reward
  (25-episode moving average, left) and training loss (right) over the 533 or 666~episodes each agent was
  trained for. The four path-scoring agents converge to within 0.01 of one another at roughly $+0.81$,
  while the flat DQN plateaus near $+0.05$ with a \emph{rising} loss despite the identical budget, data,
  and optimizer -- the evidence that its deficit is architectural rather than a matter of training time.
  The simple scoring agent's loss also diverges, which is why the enhanced variant is used throughout.}
  \label{fig:app:training}
\end{figure}

\Cref{fig:app:training} substantiates the convergence discussion of \cref{sec:p1impl:agents}. The
per-episode statistics also carry the stratified schedule's profile label for each episode, giving
roughly 111~episodes per intent for the conditional agents; these per-intent curves are the only place
the two probe-frugality profiles of \cref{tab:app:profiles} appear, since \cref{sec:p1eval:setup}
excludes them from the conditioning studies.

\begin{figure}[h]
  \centering
  \includegraphics[width=\linewidth]{p2app_training_curves.png}
  \caption[Training curves for the multipath agents]{Training curves for the multipath agents of \cref{sec:p2impl:agents}: per-episode application
  QoE, frame-loss fraction, and VMAF over the 100~episodes of each run, smoothed with a five-episode
  moving average, one line per training seed. The first three columns are the per-intent specialists of
  \cref{tab:p2results}, whose three seeds are tight enough to overplot; the fourth is the attention-pool
  arm of \cref{sec:p2eval:ablation}, whose third seed is the reason that arm's training-seed spread is
  0.151 against the mean-pool's 0.013. That seed tracks the other two for roughly 25~episodes, diverges
  into a loss excursion peaking near 0.44, and recovers only partially by episode 100 while its VMAF
  decays from 80 to below 60. A training curve is the most direct evidence available that a
  training-seed spread is a collapse rather than noise, which is why the arm is reported rather than
  dropped. Source: the \texttt{history} array of \texttt{rl-mpquic/runs/*/stats.json}.}
  \label{fig:app:p2training}
\end{figure}

\Cref{fig:app:p2training} is the multipath counterpart of \cref{fig:app:training} and substantiates two
claims from \cref{ch:multipath}. The three intent specialists converge within roughly 40~episodes and
their three seeds are visually indistinguishable thereafter, which is what a training-seed spread of
0.018--0.023 looks like; the attention-pool arm shows what 0.151 looks like, and the shape --- a clean
start, an abrupt divergence, a partial recovery --- is a collapse and not a slower learner.

\section{Multipath System: The Path-Count Sweep}
\label{app:p2pathcount}

The path-count study of \cref{sec:p2eval:pathcount} sweeps $N \in \{2, 3, 4, 6, 8, 12\}$ against
checkpoints trained at $N = 6$. \Cref{tab:app:p2pathcount} states how each configuration was built,
because the construction is what keeps the sweep a measurement of the path count rather than of the
bandwidth behind it. Everything not listed --- video, episode, reward, learner, and the four
non-stationary mechanisms with their hazard rates --- is copied unchanged from the nominal configuration
of \cref{tab:app:p2topo}, and the $N = 6$ entry of the primary sweep is config-object-identical to it.

\begin{table}[h]
  \centering
  \caption[Construction of the path-count sweep]{Construction of the path-count sweep (\cref{sec:p2eval:pathcount}); the two sweeps differ
  only in how per-path rates are set. Source:
  \texttt{rl-mpquic/\allowbreak scripts/\allowbreak make\_pathcount\_configs.py} and
  \texttt{configs/\allowbreak pathcount/*.yaml}.}
  \label{tab:app:p2pathcount}
  \small
  \begin{tabular}{@{}l p{0.58\linewidth}@{}}
    \toprule
    Choice & Value \\
    \midrule
    Role templates & the six of \cref{tab:app:p2topo} verbatim: (delay, cross-traffic fraction, nominal
                     rate) for each of the six path characters \\
    Subset for $N < 6$ & clean low-RTT, latency trap, congested, clean --- in that order, \emph{not} the
                     listing order, which would give $N = 2$ two near-identical clean paths \\
    Extension for $N > 6$ & the six roles cycle; the $k$-th repeat has its delay scaled by $1 + 0.2k$ so
                     repeated roles stay distinguishable \\
    Rates, \texttt{fixed\_total} & scaled so that every $N$ totals \SI{15.5}{\mega\bit\per\second}
                     (primary sweep) \\
    Rates, \texttt{fixed\_perpath} & nominal per path, so the aggregate is $15.5 \cdot N/6$
                     (contrast sweep) \\
    Minimum live paths & $\max(1, \lceil N/2 \rceil)$, i.e.\ 1/2/2/3/4/6, preserving the nominal
                     three-of-six fraction \\
    Shared-bottleneck groups & one per completed cycle of six: none at $N = 2, 3, 4$; one at $N = 6, 8$;
                     two at $N = 12$ \\
    \bottomrule
  \end{tabular}
\end{table}

Two consequences of the fixed-total construction bound what \cref{sec:p2eval:pathcount} claims. Per-path
rate ranges from \SI{8.45}{\mega\bit\per\second} at $N = 2$ down to \SI{1.29}{\mega\bit\per\second} at
$N = 12$, against the ${\approx}\SI{2.5}{\mega\bit\per\second}$ the checkpoints trained on; and the
minimum-live-path count cannot both preserve the nominal 50\% fraction and stay an integer at $N = 3$,
where it is 2 of 3. The absence of a shared-bottleneck group at $N = 2, 3, 4$ makes those points easier
along an axis other than the path count, which is stated in the text as well.

\section{Multipath System: The Regime Sweep}
\label{app:p2regimes}

The regime map of \cref{fig:p2eval:regimemap} varies the network's harshness and nothing else. The two
held-out configurations \texttt{dynamic\_ood\_\{mild,harsh\}.yaml} differ from the nominal one only in
the \texttt{dynamics:} block, so the axis is built by interpolating those same six parameters ---
\texttt{churn\_down\_rate}, \texttt{regime\_rate}, \texttt{regime\_lo}, \texttt{burst\_rate},
\texttt{burst\_intensity}, and \texttt{corr\_intensity} --- between them. Topology, video, reward, and
learner settings are the nominal ones throughout. The nominal configuration is not the exact midpoint of
the two held-out ones, so it is inserted as its own row rather than synthesized, and rows r1 and r3 are
the two half-way points. The row coordinate reported in \cref{fig:p2eval:regimemap}(b) is not that
parameter bundle but the \emph{even} baseline's measured mean loss fraction in each cell, which runs
0.031, 0.034, 0.039, 0.049, and 0.067 down the five rows. Source:
\texttt{rl-mpquic/\allowbreak scripts/\allowbreak make\_regime\_configs.py} and
\texttt{configs/\allowbreak regime\_r\{0..4\}\_*.yaml}.

\section{Code and Artifact Availability}
\label{app:artifacts}

Two codebases produce every number in this thesis. The single-path selector of \cref{ch:single_path} is
implemented in one repository holding the topology generator, the beaconing and traffic simulators, the
agent family of \cref{sec:p1impl:agents}, and the six-stage pipeline of \cref{sec:p1impl:pipeline}. The
multipath system of \cref{ch:multipath} is implemented in a second, \texttt{rl-mpquic}, holding the two
data-plane backends, the agent pair, the evaluation harness, and the configuration files cited
throughout that chapter.

\gap{insert the two repository URLs and the commit hash of each that produced the reported results,
together with the versions of Python, PyTorch, NS-3, and ns3-ai used. These are the only facts this
section still needs.}

Every table and figure in \cref{ch:multipath} names the run directory behind it in its caption, and
those directories live under \texttt{rl-mpquic/runs/}; \texttt{rl-mpquic/docs/RESULTS.md} maps each
reported result to the command that produced it, and the single-path pipeline writes the equivalent
artifacts into one timestamped run directory per invocation (\cref{sec:p1impl:pipeline}). Three limits
on reproducibility belong with that. The trained checkpoints are not under version control, so results
can be regenerated by retraining but not by replaying our exact weights. The NS-3 arms additionally
require a built simulator with the ns3-ai shared-memory bridge, and cannot be run concurrently, for the
reason given in \cref{sec:p2eval:ns3native}. And because the spread over training seeds exceeds the
evaluation-seed interval by an order of magnitude on this problem (\cref{sec:p2eval:seeds}), a
reproduction should be read against the reported training-seed spreads rather than against the reported
means.
